\chapter*{Заключение}
\addcontentsline{toc}{chapter}{Заключение}
Была разработана система ролевой игры с виртуальным мастером, обеспечивающая структурированное управление игровым процессом, контроль действий игрока и реалистичное поведение неигровых персонажей на основе их эмоционального состояния и намерений.

\section*{Аналитические результаты}

Выполнен анализ существующих коммерческих платформ (AI Dungeon, AI Dungeon-Chi, Dungeon Master AI, HyperWrite AI) и исследовательских работ в области применения больших языковых моделей для ролевых игр. Выявлены ключевые проблемы существующих систем: чрезмерная уступчивость ИИ, отсутствие слежения за состоянием мира, неспособность к отказу или корректирующей реакции. Установлено, что ни одна из существующих систем не обеспечивает необходимый уровень контроля действий игрока, структурированного представления состояния мира, интеграции механических правил и управления неигровыми персонажами.

\section*{Теоретические результаты}

Разработана формальная модель системы ролевой игры, включающая иерархическую структуру компонентов и формализацию понятий действия, директивы и повествования. Разработана модель управления эмоциональным состоянием персонажей на основе моральных схем, использующая векторные представления эмоций и намерений. Разработаны алгоритмы работы основных компонентов системы: игрового цикла, выдачи директив, разрешения действий и генерации ответов персонажей. Предложена обобщенная архитектура системы, организованная в виде иерархии подсистем с четко определенными интерфейсами взаимодействия.

\section*{Инженерные результаты}

Спроектирована модульная архитектура системы, включающая подсистемы управления игровым циклом, управления сценами, выдачи директив, управления персонажами, эмоционального моделирования, разрешения действий и взаимодействия с большими языковыми моделями. Разработаны интерфейсы взаимодействия между компонентами, реализованные через асинхронные методы. Интегрирована система управления эмоциональным состоянием персонажей, поддерживающая одиночные и множественные моральные схемы. Разработана архитектура компонента выдачи директив и системы разрешения действий, интегрирующей механические элементы игры с генерацией повествования.

Также есть смысл привести предполагаемые направления для будущей работы.

Реализована полнофункциональная система ролевой игры с виртуальным мастером на языке Python, представляющая собой консольное приложение. Проведено тестирование системы, основываясь на обработке краевых случаев. Проведено сравнение с существующими коммерческими аналогами: система превосходит их по контролю действий игрока, структурированному представлению состояния мира, интеграции механических правил, управлению персонажами с учетом эмоционального состояния и проверке корректности действий. Система демонстрирует более структурированный и контролируемый игровой процесс по сравнению с существующими генеративными решениями, сохраняя при этом гибкость и вариативность повествования.

%%% Local Variables:
%%% TeX-engine: xetex
%%% eval: (setq-local TeX-master (concat "../" (seq-find (-cut string-match ".*-3-pz\.tex$" <>) (directory-files ".."))))
%%% End:
